\documentclass[12pt]{article}
\usepackage[utf8]{inputenc}
\usepackage[T1]{fontenc}
\usepackage{hyperref}
\title{Yindjilandji Language Resources for K-12 Curriculum Development}
\author{Compiled by Claude Sonnet 5.5 (claude-sonnet-5-5) for Andr\'as Kornai}
\date{}
\begin{document}
\maketitle

\section{Introduction}
Yindjilandji (also spelled Indjilandji, Indjiladji, Indjilindji, Injeelanjee; ISO 639-3 yil; Glottocode yind1248) is an extinct Australian language. According to the sources found, its country lay on the Northern Territory--Queensland border region near the headwaters of the Gregory River and Lawn Hill Creek. Linguists usually treat it as a southern Ngarna-group variety, and some treat it as a dialect of Wagawa. There are no known native speakers, and I found no school materials. Any curriculum would have to be built from archival documentation, in consultation with the descendant community.

\section{Primary Educational Resources}
No primers, textbooks or school programmes for Yindjilandji turned up in the search.

\section{Linguistic and Reference Materials}
\begin{itemize}
\item Wikipedia, ``Yindjilandji'': overview of location, classification and alternative names. \url{https://en.wikipedia.org/wiki/Yindjilandji}
\item AIATSIS Austlang entry G14, Indjilandji: spelling variants, and notes on an 1891 Eric Henderson vocabulary from near Yelvertoft Station, Queensland. \url{https://collection.aiatsis.gov.au/austlang/language/g14}
\item Gavan Breen Yindjilandji Collection (Research Data Australia): record of fieldwork material. \url{https://researchdata.edu.au/gavan-breen-yindjilandji-collection/967828}
\end{itemize}

\section{Religious and Community Resources}
None found.

\section{Digital and Audio Resources}
\begin{itemize}
\item PARADISEC catalogue, collection GB34, item 003: 1979 Gavan Breen recordings (vocabulary, elicitation and glossed sentences recorded at Gregory River, Queensland, with speaker Ivy George). Access conditions were not checked. \url{https://catalog.nabu-prod.paradisec.org.au/collections/GB34/items/003}
\item University of Sydney Archives search records (nodes 19408 and 19409), which appeared in results for this language; their content was not examined. \url{https://archives-search.sydney.edu.au/nodes/view/19408} and \url{https://archives-search.sydney.edu.au/nodes/view/19409}
\end{itemize}

\section{Curriculum Development Context}
Very little exists beyond archival word lists and the Breen recordings. A heritage or revitalization curriculum would be small in scope: awareness of the language and its country, plus whatever vocabulary the archives support. It should be developed only with the involvement of the relevant Aboriginal community and with respect for archive access protocols. Because the spellings vary so much, searches should try all of the name variants above.
\end{document}